\documentclass[10pt,a4paper]{amsart}
\usepackage{amsmath,amssymb,amsthm,graphicx}
\usepackage{verbatim}
\newenvironment{reference}{\comment}{\endcomment}
%\newenvironment{reference}{}{}
\newenvironment{slogan}{\comment}{\endcomment}
\newenvironment{history}{\comment}{\endcomment}
% For commutative diagrams we use Xy-pic
\usepackage[all]{xy}

% We use 2cell for 2-commutative diagrams.
\xyoption{2cell}
% Theorem environments.
%
\theoremstyle{plain}
\newtheorem{theorem}[subsection]{Theorem}
\newtheorem{proposition}[subsection]{Proposition}
\newtheorem{lemma}[subsection]{Lemma}

\theoremstyle{definition}
\newtheorem{definition}[subsection]{Definition}
\newtheorem{example}[subsection]{Example}
\newtheorem{exercise}[subsection]{Exercise}
\newtheorem{situation}[subsection]{Situation}

\theoremstyle{remark}
\newtheorem{remark}[subsection]{Remark}
\newtheorem{remarks}[subsection]{Remarks}

\numberwithin{equation}{subsection}

% Macros
%
\def\lim{\mathop{\mathrm{lim}}\nolimits}
\def\colim{\mathop{\mathrm{colim}}\nolimits}
\def\Spec{\mathop{\mathrm{Spec}}}
\def\Hom{\mathop{\mathrm{Hom}}\nolimits}
\def\Ext{\mathop{\mathrm{Ext}}\nolimits}
\def\SheafHom{\mathop{\mathcal{H}\!\mathit{om}}\nolimits}
\def\SheafExt{\mathop{\mathcal{E}\!\mathit{xt}}\nolimits}
\def\Sch{\mathit{Sch}}
\def\Mor{\mathop{\mathrm{Mor}}\nolimits}
\def\Ob{\mathop{\mathrm{Ob}}\nolimits}
\def\Sh{\mathop{\mathit{Sh}}\nolimits}
\def\NL{\mathop{N\!L}\nolimits}
\def\CH{\mathop{\mathrm{CH}}\nolimits}
\def\proetale{{pro\text{-}\acute{e}tale}}
\def\etale{{\acute{e}tale}}
\def\QCoh{\mathit{QCoh}}
\def\Ker{\mathop{\mathrm{Ker}}}
\def\Im{\mathop{\mathrm{Im}}}
\def\Coker{\mathop{\mathrm{Coker}}}
\def\Coim{\mathop{\mathrm{Coim}}}

% Boxtimes
%
\DeclareMathSymbol{\boxtimes}{\mathbin}{AMSa}{"02}

%
% Macros for moduli stacks/spaces
%
\def\QCohstack{\mathcal{QC}\!\mathit{oh}}
\def\Cohstack{\mathcal{C}\!\mathit{oh}}
\def\Spacesstack{\mathcal{S}\!\mathit{paces}}
\def\Quotfunctor{\mathrm{Quot}}
\def\Hilbfunctor{\mathrm{Hilb}}
\def\Curvesstack{\mathcal{C}\!\mathit{urves}}
\def\Polarizedstack{\mathcal{P}\!\mathit{olarized}}
\def\Complexesstack{\mathcal{C}\!\mathit{omplexes}}
% \Pic is the operator that assigns to X its picard group, usage \Pic(X)
% \Picardstack_{X/B} denotes the Picard stack of X over B
% \Picardfunctor_{X/B} denotes the Picard functor of X over B
\def\Pic{\mathop{\mathrm{Pic}}\nolimits}
\def\Picardstack{\mathcal{P}\!\mathit{ic}}
\def\Picardfunctor{\mathrm{Pic}}
\def\Deformationcategory{\mathcal{D}\!\mathit{ef}}
\usepackage{cleveref}
\crefname{originaldefinition}{original definition}{original definitions}
\Crefname{originaldefinition}{Original definition}{Original definitions}
\makeatletter
\newlabel{definition-gysin-homomorphism}{{02T8}{source}{}{sourcepointer-1}{}}
\newlabel{definition-gysin-homomorphism@cref}{{[originaldefinition][0][]02T8}{[1][0][]source}{}{}{}}
\makeatother
\crefname{originallemma}{original lemma}{original lemmas}
\Crefname{originallemma}{Original lemma}{Original lemmas}
\makeatletter
\newlabel{lemma-decompose-section}{{0F97}{source}{}{sourcepointer-2}{}}
\newlabel{lemma-decompose-section@cref}{{[originallemma][0][]0F97}{[1][0][]source}{}{}{}}
\makeatother
\crefname{originallemma}{original lemma}{original lemmas}
\Crefname{originallemma}{Original lemma}{Original lemmas}
\makeatletter
\newlabel{lemma-decompose-section-formulae}{{0F98}{source}{}{sourcepointer-3}{}}
\newlabel{lemma-decompose-section-formulae@cref}{{[originallemma][0][]0F98}{[1][0][]source}{}{}{}}
\makeatother
\crefname{originallemma}{original lemma}{original lemmas}
\Crefname{originallemma}{Original lemma}{Original lemmas}
\makeatletter
\newlabel{lemma-linebundle}{{0B74}{source}{}{sourcepointer-4}{}}
\newlabel{lemma-linebundle@cref}{{[originallemma][0][]0B74}{[1][0][]source}{}{}{}}
\makeatother
\crefname{originalreference}{original reference}{original references}
\Crefname{originalreference}{Original reference}{Original references}
\makeatletter
\newlabel{situation-setup}{{02QL}{source}{}{sourcepointer-5}{}}
\newlabel{situation-setup@cref}{{[originalreference][0][]02QL}{[1][0][]source}{}{}{}}
\makeatother
\title{Deforming a triple so that the divisor is cut out by one global function, with finite type flat relative dimension one maps and injective pullback on Chow groups}
\author{The Stacks Project Authors}
\date{stacks-project:chow.tex:local-snapshot-20261012; GNU FDL 1.2 or later}
\begin{document}
\maketitle
\noindent\textit{Editorial note.} All delivered mathematical text is the author's own, copied from the pinned Stacks Project chapter file chow.tex; no statement or proof step was rewritten, reordered or removed.\par\medskip
\noindent\textit{Editorial note.} Preamble provenance: the pinned chapter file opens with \textbackslash{}input\{preamble\} and carries no preamble text of its own, so the project's single Stacks-wide preamble staged at sources/stacks-chow/preamble.tex is excerpted verbatim as three spans, lines 8-12 (the verbatim package together with the project's own reference, slogan and history display-environment definitions, including the author's commented-out older form of the reference environment), lines 14-18 (Xy-pic setup) and lines 154-217 (the native theorem declarations and mathematical macros). The lines 8-12 span is required because a counted statement or its surrounding text may use those display environments, which the project typesets as comments. Nothing was added, removed or redefined; the wrapper class and the amsmath/amssymb/amsthm/graphicx base packages are supplied by the assembler.\par\medskip
\noindent\textit{Editorial note.} Same-chapter pointers are located at their \textbackslash{}label lines inside the pinned chapter file; cross-chapter pointers are located in the project's own tag list staged in the snapshot (sources/stacks-chow/tags.txt), at the verbatim lines carrying their canonical tags.\par\medskip
\noindent\textit{Allowed standard prerequisites (editorial).}\begin{itemize}
\item The Gysin homomorphism attached to a triple, which makes the statement meaningful for the deformed triple. Same chapter chow.tex, Definition \textbackslash{}ref\{definition-gysin-homomorphism\}, invoked at proof line 5965. The deformed data \$(\textbackslash{}mathcal\{O\}\_\{X'\}, s', i' : D' \textbackslash{}to X')\$ is again a triple as in the definition, so the pullback \$(i')\textasciicircum{}*\$ is defined and the identity in item (4) makes sense.
\item Decomposition of a global section of a tensor product of line bundles into two sections. Same chapter chow.tex, label lemma-decompose-section (invoked at proof line 5970). Applied with \$\textbackslash{}mathcal\{L\}\_1 = \textbackslash{}mathcal\{O\}\_X\$ and \$\textbackslash{}mathcal\{L\}\_2 = \textbackslash{}mathcal\{L\}\$ to the section \$s\$, it produces the commutative diagram together with the finite type, flat of relative dimension one morphisms \$p\$ and \$g\$ and the presentation of \$D' = D\_1\$ as the zero scheme of a global section.
\item The formulae attached to that decomposition. Same chapter chow.tex, label lemma-decompose-section-formulae (invoked at proof line 5973). Supplies the equality \$p\textasciicircum{}*i\textasciicircum{}* = (i')\textasciicircum{}*g\textasciicircum{}*\$ between the maps on Chow groups.
\item Injectivity of the pullback along the projection onto the zero scheme in the decomposed situation. Same chapter chow.tex, label lemma-linebundle (invoked at proof line 5975). Gives \$p\textasciicircum{}* : \textbackslash{}CH\_k(D) \textbackslash{}to \textbackslash{}CH\_\{k+1\}(D')\$ injective for all \$k\$.
\end{itemize}\medskip
\ifdefined\hypertarget\hypertarget{sourcepointer-1}{}\fi
\noindent\textit{Source reference.} Original definition 02T8: sources/stacks-chow/pinned-original-source.tex, source lines 5069--5069. Navigation pointer for \textbackslash{}ref\{definition-gysin-homomorphism\} used inside the counted span (line 5943); the argument that consumes it is retained in the span.\par\medskip
\ifdefined\hypertarget\hypertarget{sourcepointer-2}{}\fi
\noindent\textit{Source reference.} Original lemma 0F97: sources/stacks-chow/pinned-original-source.tex, source lines 5822--5822. Navigation pointer for \textbackslash{}ref\{lemma-decompose-section\} used inside the counted span (line 5970); the argument that consumes it is retained in the span.\par\medskip
\ifdefined\hypertarget\hypertarget{sourcepointer-3}{}\fi
\noindent\textit{Source reference.} Original lemma 0F98: sources/stacks-chow/pinned-original-source.tex, source lines 5884--5884. Navigation pointer for \textbackslash{}ref\{lemma-decompose-section-formulae\} used inside the counted span (line 5973); the argument that consumes it is retained in the span.\par\medskip
\ifdefined\hypertarget\hypertarget{sourcepointer-4}{}\fi
\noindent\textit{Source reference.} Original lemma 0B74: sources/stacks-chow/pinned-original-source.tex, source lines 5737--5737. Navigation pointer for \textbackslash{}ref\{lemma-linebundle\} used inside the counted span (line 5975); the argument that consumes it is retained in the span.\par\medskip
\ifdefined\hypertarget\hypertarget{sourcepointer-5}{}\fi
\noindent\textit{Source reference.} Original reference 02QL: sources/stacks-chow/pinned-original-source.tex, source lines 1372--1372. Navigation pointer for \textbackslash{}ref\{situation-setup\} used inside the counted span (line 5941); the argument that consumes it is retained in the span.\par\medskip
\begin{lemma}
\label{lemma-normal-cone-effective-Cartier}
In Situation \ref{situation-setup} let $X$ be a scheme locally
of finite type over $S$. Let $(\mathcal{L}, s, i : D \to X)$
be a triple as in Definition \ref{definition-gysin-homomorphism}.
There exists a commutative diagram
$$
\xymatrix{
D' \ar[r]_{i'} \ar[d]_p & X' \ar[d]^g \\
D \ar[r]^i & X
}
$$
such that
\begin{enumerate}
\item $p$ and $g$ are of finite type and flat of relative dimension $1$,
\item $p^* : \CH_k(D) \to \CH_{k + 1}(D')$ is injective for all $k$,
\item $D' \subset X'$ is the zero scheme of a global section
$s' \in \Gamma(X', \mathcal{O}_{X'})$,
\item $p^*i^* = (i')^*g^*$ as maps $\CH_k(X) \to \CH_k(D')$.
\end{enumerate}
Moreover, these properties remain true after arbitrary base change
by morphisms $Y \to X$ which are locally of finite type.
\end{lemma}

\begin{proof}
Observe that $(i')^*$ is defined because we have the triple
$(\mathcal{O}_{X'}, s', i' : D' \to X')$ as in
Definition \ref{definition-gysin-homomorphism}. Thus the statement makes sense.

\medskip\noindent
Set $\mathcal{L}_1 = \mathcal{O}_X$, $\mathcal{L}_2 = \mathcal{L}$
and apply Lemma \ref{lemma-decompose-section} with the section $s$ of
$\mathcal{L} = \mathcal{L}_1 \otimes_{\mathcal{O}_X} \mathcal{L}_2$.
Take $D' = D_1$. The results now follow from the lemma, from
Lemma \ref{lemma-decompose-section-formulae}
and injectivity by
Lemma \ref{lemma-linebundle}.
\end{proof}

\end{document}
